

\paragraph{Keep carried tax payments separate from the new assessment.}
The first current-tax assessment after carry has a different opening from the
same-period reassessment above. Current-period expense begins at zero, but the
book can still owe tax from the preceding period. Paying that opening liability
settles an old assessment; it does not reduce the new period's tax expense.
Calculating paid tax as current expense less the outstanding payable would
therefore be wrong for this opening.

The request binds the complete assessed-carry execution and begins its tax
interval at the new financial opening and loss cursor. It requires a completed
incoming execution and no previous assessment in the new period. Before that
first assessment, the selected payment kernel can consume only the carried
liability. Each payment can therefore be attributed to the prior assessment
from its exact event identity, date, amount and cash-flow record. A stopped
incoming execution is retained for inspection but is not silently assessed.

Let $P_0$ denote the opening tax payable, $S$ the completed payments against
that carried liability and $C$ the new conditional assessment. The two liability
records reconcile to the financial book as
\[
 P_{\mathrm{carried}}=P_0-S,\qquad
 P_{\mathrm{new}}=C,\qquad
 P_{\mathrm{total}}=P_0-S+C,\qquad 0\leq S\leq P_0.
\]
The first record retains the prior assessment's interval and identity with its
new-period settlements. The second records the new interval, assessment and
initially empty settlement history. The journal debits current-tax expense by
$C$ and credits the payable by the same amount. Cash, deferred balances and the
old settlements remain unchanged. Even a fully paid old liability keeps its
history rather than disappearing from the explanation.

For opening payable $19/4$ and new tax of 5, paying all $19/4$ leaves a total
payable of 5. Paying only 1 leaves $35/4$, and paying nothing leaves $39/4$.
The new expense is 5 in all three cases. In the fully paid example, after the
new-period depreciation of 10, recorded profit is $-15$, cash is $1423/4$,
assets are $2735/4$, liabilities are 115 and equity is $2275/4$. The increase
in the new payable cancels its expense in the indirect cash calculation; the
old payment remains the operating cash outflow of $19/4$.

Taxable profit is calculated from the new period's pretax accounts and explicit
adjustments. Each native return adjustment must appear exactly once, and the
complete tax-base declarations must reconcile at account and lease-right level.
Those declarations do not themselves assess deferred tax. In a return example,
new-period recorded profit of 6 and its native adjustment of $-24$ produce
taxable loss of 18. The method retains a new loss quantity of 18, available at
the interval's exclusive endpoint, without a refund or an automatic deferred
asset.

Loss use starts from the actual carried inventory. If a loss of 19 remains
available, explicit use against taxable profit of 20 leaves taxable profit of
1 and tax of $1/4$ at the declared rate. The history continues from the prior
cursor, with total use of 19. The selected close-day rule separately checks
expiry: a loss expiring before the usage day cannot be used, while a use before
expiry consumes only the declared quantity. Neither the previous deferred-tax
recognition nor its recovery capacity consumes or replenishes the loss inventory.

Applying the proposal compares the full predecessor, tax calculations, loss
movements, both liability records and financial statements. The checked scope
is a first assessment without credits, refunds or prepayments. Once both tax
vintages exist, subsequent payments need explicit allocation; the original
aggregate payment interface no longer supplies enough information. New-period
deferred assessment, repeated current assessments and later carry require their
own compatible methods. Source admission and complete annual event coverage
remain unresolved.
